\section{Conjugacy classes in fibre products}\label{sec:conjugacy}

To compute Gassmann classes of subgroups of the finite quotients of
\cref{sec:quotients} one needs their conjugacy classes. By the fibre
product decomposition (\cref{thm:fibre-product}) these come from the
prime power pieces, but not simply as products.

\begin{setup}[Fibre products over an abelian group]\label{setup:fibre-product}
Let $A$ be an abelian group and $\pi_i \colon G_i \to A$, $i = 1, \ldots, r$,
surjective homomorphisms of groups with kernels $K_i$. Let
$G := \{(g_i)_i \in \prod_i G_i : \pi_1(g_1) = \cdots = \pi_r(g_r)\}$ and
$\pi \colon G \to A$ the common value. For $a \in A$ write
$\Fibre{G_i}{a} := \pi_i^{-1}(a)$ and $\Fibre{G}{a} := \pi^{-1}(a)$. For
$x_i \in G_i$ put $\CentralizerImage{x_i} := \pi_i(\Centralizer{G_i}{x_i})
\subseteq A$.
\end{setup}

\begin{lemma}[Types are conjugation invariant]\label{lem:types-invariant}
In \cref{setup:fibre-product}, conjugation in $G$ (respectively $G_i$)
preserves each $\Fibre{G}{a}$ (respectively $\Fibre{G_i}{a}$).
\end{lemma}

\begin{proof}
$A$ is abelian.
\end{proof}

\begin{lemma}[Conjugation by the kernel is componentwise]
\label{lem:kernel-componentwise}
In \cref{setup:fibre-product}, $K := \prod_i K_i$ is a normal subgroup of
$G$ with $G/K \cong A$, and the $K$-conjugacy classes in $\Fibre{G}{a}$ are
the products of $K_i$-conjugacy classes in the $\Fibre{G_i}{a}$.
\end{lemma}

\begin{proof}
$K$ is the kernel of $\pi$, which is surjective; and $K$ acts on
$\prod_i G_i$ componentwise.
\end{proof}

\begin{lemma}[The diagonal action]\label{lem:diagonal-action}
\slogan{Conjugacy classes of a fibre product are orbits of the base group
on tuples of local kernel classes.}
In \cref{setup:fibre-product}, for $b \in A$ and a $K_i$-conjugacy class
$C \subseteq \Fibre{G_i}{a}$ put $b\FibreAction C := hCh^{-1}$ for any
$h \in \Fibre{G_i}{b}$. This is a well defined action of $A$ on the set
$\ConjClassesBy{K_i}{\Fibre{G_i}{a}}$ of $K_i$-classes, the stabiliser of
the class of $x_i$ is $\CentralizerImage{x_i}$, and
\[
  \ConjClassesBy{G}{\Fibre{G}{a}} \cong
  \Bigl(\prod_i \ConjClassesBy{K_i}{\Fibre{G_i}{a}}\Bigr)/A,
\]
with $A$ acting diagonally.
\end{lemma}

\begin{proof}
Two elements of $\Fibre{G_i}{b}$ differ by an element of $K_i$. The
stabiliser statement: $hx_ih^{-1}$ is $K_i$-conjugate to $x_i$ if and only
if $kh$ centralises $x_i$ for some $k \in K_i$. The last statement follows
from \cref{lem:kernel-componentwise}, since every element of $G$ is an
element of $K$ times $(h_i)_i$ with all $h_i$ of the same type.
\end{proof}

\begin{proposition}[Fibres over local classes]\label{prop:fibres-over-local}
\slogan{Over a tuple of local classes, the global classes form a torsor
under $\Coker(A \to \prod A/S_i)$.}
In \cref{setup:fibre-product}, the map
$\ConjClassesBy{G}{\Fibre{G}{a}} \to \prod_i \ConjClassesBy{G_i}{\Fibre{G_i}{a}}$
is surjective. Its fibre over $(\ConjClassOf{x_i})_i$ is a torsor under
\[
  \Coker\Bigl(A \xrightarrow{\ \mathrm{diag}\ } \prod_i A/\CentralizerImage{x_i}\Bigr),
\]
of cardinality $\prod_i\Index{A}{\CentralizerImage{x_i}}\big/
\Index{A}{\bigcap_i\CentralizerImage{x_i}}$. Concretely, $x = (x_i)_i$ and
$y = (y_i)_i$ in $\Fibre{G}{a}$ are conjugate in $G$ if and only if there
are $u_i \in G_i$ with $y_i = u_ix_iu_i^{-1}$ such that the element
$(\pi_i(u_i) \bmod \CentralizerImage{x_i})_i$ lies in the image of the
diagonal; this element does not depend on the choice of the $u_i$ modulo
the image of the diagonal.
\end{proposition}

\begin{proof}
The $G_i$-class of $x_i$ is the $A$-orbit of its $K_i$-class
(\cref{lem:diagonal-action}), isomorphic to $A/\CentralizerImage{x_i}$ as
an $A$-set; so the fibre is the quotient of $\prod_i A/\CentralizerImage{x_i}$
by the diagonal action of $A$, which since $A$ is abelian is the cokernel
of the diagonal map, acting simply transitively. The cardinality is
$\#\prod_i(A/\CentralizerImage{x_i})$ divided by the order of the image of
$A$, which is $\Index{A}{\bigcap_i\CentralizerImage{x_i}}$.
\end{proof}

\begin{remark}[Canonical coordinates]\label{rem:basepoints}
The group $\Coker(A \to \prod_i A/\CentralizerImage{x_i})$ is canonical,
but to identify the fibre with it one has to choose a $K_i$-class inside
each $G_i$-class (a basepoint). If the local conjugacy classes are labelled
by invariants that are fine enough to distinguish $K_i$-classes, one can
take the basepoint with the smallest label, which makes the coordinates
canonical.
\end{remark}

\begin{example}[Application to the finite quotients]
\label{ex:conjugacy-quotients}
By the fibre product decomposition (\cref{thm:fibre-product}), the
conjugacy classes of $\AdelicQuotientGroup{N}{M}$ and
$\LiftedQuotientGroup{N}{M}$ are described by \cref{prop:fibres-over-local}
with $A = \TypeGroup{N}$ and $G_i$ the prime power pieces. The kernels
$K_i$ are the groups of \cref{prop:det-one-quotients} for $M = p_i^{k_i}$.
For $p_i \nmid N$ they are (determinant one parts of)
$\GL{2}{\ResidueRing{p_i^{k_i}}}$, and a type $e$ acts on their classes by
conjugation by a matrix of determinant $e$.
\end{example}
